
\documentclass[../../main.tex]{subfiles}
% \documentclass[a4paper,10pt]{ltjsarticle}
\begin{document}
% \documentclass[a4paper,12pt]{article}
% \usepackage[utf8]{inputenc}
% \usepackage[japanese]{babel}
% \usepackage{xeCJK}
% \usepackage{geometry}
% \usepackage{array}
% \usepackage{booktabs}
% \usepackage{fancybox}

% % ページ設定
% \geometry{top=25mm, bottom=25mm, left=20mm, right=20mm}
% \pagestyle{empty}

% % フォント設定
% \setCJKmainfont{Noto Sans CJK JP}

% \begin{document}

% 年度の大きな表示
\begin{center}
  {\fontsize{48pt}{58pt}\selectfont \textbf{1998}\normalsize\textbf{年度}}
\end{center}

\vspace{10mm}

% 本文


\vspace{15mm}

% 出題テーマと難易度の表
\noindent\textbf{出題テーマと難易度}

\vspace{5mm}

\begin{center}
  \shadowbox{%
    \renewcommand{\arraystretch}{1.6}
    \begin{tabular}{|c|p{8cm}|c|}
      \hline
      \textbf{問題番号} & \textbf{テーマ} & \textbf{難易度} \\
      \hline
      第1問           & 極限           & 易            \\
      \hline
      第2問           & 空間図形         & 普            \\
      \hline
    \end{tabular}%
    \renewcommand{\arraystretch}{1.0}
  }
\end{center}

\vspace{15mm}

% 解答
\noindent\textbf{解答}

\vspace{5mm}

\begin{center}
  \shadowbox{%
    \renewcommand{\arraystretch}{1.6}
    \begin{tabular}{|c|c|p{8cm}|}
      \hline
      \textbf{問題番号}        & \textbf{小問} & \textbf{答え}                                         \\
      \hline
      第1問                  & -           & $\displaystyle     \lim_{n\to\infty} x_n =
      \begin{dcases}
          0                     & (a+b>1)          \\
          \log 2                & (b=0, a=1)       \\
          \frac{2^{1-a}-1}{1-a} & (a+b=1, a \ne 1)
        \end{dcases}$                           \\
      \hline
      \multirow{3}{*}{第2問} & (1)         & $\displaystyle 2x^2+y^2=2$                          \\
      \cline{2-3}
                           & (2)         & $\displaystyle  R =
      \begin{dcases}
          1-z^2 \le y^2 \le (z+\sqrt{2})^2          & 0 \le z \le \frac{\sqrt{2}}{2} \\
          (z-\sqrt{2})^2 \le y^2 \le (z+\sqrt{2})^2 & \frac{\sqrt{2}}{2} \le z \le 1
        \end{dcases} $ \\
      \cline{2-3}
                           & (3)         & $\displaystyle V = \frac{49}{3}\sqrt{2}\pi$         \\
      \hline
    \end{tabular}%
    \renewcommand{\arraystretch}{1.0}
  }
\end{center}

% 解答時間
\noindent\textbf{試験形態}

\vspace{5mm}

この年から試験形態が変更となり，2問で90分，つまり1問あたり45分かけられるようになった．
\begin{center}
  \shadowbox{%
    \begin{tabular}{p{8cm}}
      2問90分
    \end{tabular}%
  }
\end{center}


\end{document}